\section{What Makes AI-Supported Engineering Difficult?}
\label{sec:challenges}

The use of generative AI in engineering differs from unconstrained content generation.
Engineering information contributes to technical decisions, system descriptions, models,
requirements, and other artifacts that may remain relevant throughout the system lifecycle.
AI-supported processing must therefore consider not only whether technically plausible
content can be generated, but also how the underlying information is interpreted, which
results may be used for subsequent engineering activities, and how the resulting states
remain verifiable and traceable.

\subsection{Engineering Task and Information Basis}
\label{sec:challenge_information_basis}

The introduction of AI support first requires a sufficiently bounded engineering task.
The intended result determines which source information is relevant, which interpretation
steps are required, and which criteria can later be used to assess the result. This becomes
particularly important in brownfield environments, where relevant system information may
already exist across specifications, tables, technical documents, models, or further data
sources \cite{Maheshwari2018,Henderson2024}.

Such information does not generally form a uniform or fully formalized basis. Sources may
differ in format, degree of structure, terminology, granularity, level of abstraction, and
currency. Several sources may describe the same system aspect from different perspectives,
complement each other, or contain diverging statements. The processing task therefore
cannot be reduced to making existing documents technically accessible.

A controlled AI-supported workflow also has to distinguish the role that available
information plays within the engineering task. In the underlying investigation,
\emph{Engineering Sources} formed the domain-specific information basis from which
statements about the system could be derived. \emph{Context Sources} provided additional
information for interpretation and classification without themselves constituting the
engineering evidence for the system under consideration. The \emph{Processing and
Orchestration Context}, including prompts, agent roles, and processing configurations,
influenced how information was processed but likewise did not establish engineering facts.

This distinction is relevant beyond the investigated transformation task. Information that
supports interpretation should not implicitly acquire the same authority as information
that constitutes the engineering basis itself. At the same time, the conditions under which
a result was generated remain relevant for understanding its origin. Provenance therefore
concerns both the source information and the processing context that contributed to the
result \cite{MoreauMissier2013,Simmhan2005}.

\subsection{Interpretation, Heterogeneity, and Uncertainty}
\label{sec:challenge_interpretation}

For formally defined source and target models, transformation can be governed by explicit
rules and metamodels. Heterogeneous legacy information introduces an additional problem:
before information can be transformed into a target structure, its engineering meaning has
to be identified and placed into the relevant system context.

The heterogeneity of engineering information is therefore not only a technical format
problem but also a semantic one. Information may only become understandable in relation to
its domain context, terminology, related artifacts, or existing engineering knowledge
\cite{Bressan.2000}. This
interpretation task is one area in which probabilistic AI can provide value, because it can
generate possible interpretations of weakly structured or natural-language information.

However, probabilistic interpretation introduces uncertainty. A linguistically convincing
or formally valid result does not provide direct evidence of engineering correctness.
Studies of LLM-generated systems-engineering artifacts have shown that outputs resembling
expert artifacts can still contain technically relevant errors or unsupported assumptions
\cite{Topcu2025}. Related approaches to model generation therefore combine generative
processing with intermediate representations, additional constraints, or validation
mechanisms \cite{Cibrian.2025,Stein2026,Pan.2025}.

The underlying investigation addressed this problem by applying several configured
perspectives to the same source-grounded engineering subject. Agreement between
interpretations provided information about \emph{Consensus}, while relevant differences
were represented as \emph{Variance}. A situation in which several perspectives identified
the same unresolved information or ambiguity was represented as \emph{Shared Ambiguity}.
Such multi-perspective processing can provide additional evidence for subsequent evaluation,
but agreement alone does not establish that an interpretation is correct
\cite{Du2024,TriemDing2024}.

The more general implication is that uncertainty should remain visible within an
AI-supported engineering process. Missing information, conflicting interpretations, or
insufficient evidence should not automatically be converted into apparently complete
engineering information by generative processing.

\subsection{Engineering Authority and Controlled Processing}
\label{sec:challenge_authority}

Probabilistic interpretation and engineering decision-making represent different functions.
AI-generated interpretations can provide decision support, but the decision whether a
particular interpretation may be used for subsequent engineering activities has a different
status.

Human-in-the-loop approaches describe different forms of human participation, ranging from
supervision and feedback to explicit evaluation and decision-making
\cite{MosqueiraRey2023,Lazaros2026}. In the investigated architecture, human involvement
was therefore represented as explicit decision authority at defined transitions of the
engineering process.

The \emph{Human Engineering Authority} authorized the domain-specific use of interpreted
engineering information. This decision referred to a defined information state rather than
to an abstract source or an unrestricted AI output. A materially changed state therefore
required a new evaluation.

A separate \emph{Model Authority} addressed a different decision object: how already
authorized engineering information should be structurally placed and represented in the
target model. This distinction emerged from the observation that engineering meaning,
structural representation, and target-specific formulation are not necessarily identical
decisions.

The exact authority roles depend on the engineering task. The general challenge is to avoid
allowing probabilistic processing to implicitly determine consequential engineering states.
Once the required engineering meaning, authority, and representation rules have been
resolved, subsequent processing can increasingly rely on explicit rules and constraints.
Conversely, if the required evidence, authorization, or transformation rule is missing, the
unresolved state should remain visible rather than being silently completed by the AI
system.

\subsection{Traceability, Assurance, and Integration}
\label{sec:challenge_assurance}

AI-supported engineering creates a chain of information states rather than a single
generation event. Source information may be identified, interpreted, evaluated, authorized,
transformed, and finally represented in a technical artifact. If only the final artifact is
retained, the relationship between the result and the decisions that produced it is lost.

Provenance describes the origin and derivation of information
\cite{MoreauMissier2013,Simmhan2005}. Traceability connects the resulting information,
decision, and artifact states \cite{ISO29148}. In the underlying investigation, source-grounded evidence,
AI-generated interpretations, human authorization decisions, intermediate model states,
and the final artifact remained explicitly connected. This allowed the processing path to
be reconstructed from source to result and, conversely, from a resulting model element back
to its underlying information and decisions.

This becomes increasingly relevant when several Engineering Sources contribute to a common
engineering result. Related information may need to be considered jointly while the
individual origin and authorization state of each contribution remain distinguishable.

Finally, generation and the subsequent assessment of the resulting artifact address
different questions. The fact that an AI system can produce an artifact does not demonstrate
that the artifact satisfies the relevant engineering, formal, or technical criteria.
Appropriate assessment mechanisms therefore have to be defined separately from the
generation step \cite{Topcu2025,Cibrian.2025,Stein2026}.

Moving from a bounded AI-supported task toward productive engineering use adds an
organizational dimension. The workflow has to be integrated into existing responsibilities,
quality mechanisms, and engineering processes. In addition, increasing processing volume
also increases the need to consider the scalability of human review and authorization.
These questions extend beyond model capability and make productive AI-supported engineering
a combined information, process, and systems-engineering challenge.
